\documentclass[11pt]{article}
\usepackage[T1]{fontenc}
\usepackage{lmodern}
\usepackage[margin=1in]{geometry}
\usepackage{amsmath,amssymb,amsthm}
\usepackage{graphicx}
\usepackage[hidelinks]{hyperref}
\emergencystretch=2em

\newtheorem{theorem}{Theorem}
\newtheorem{corollary}[theorem]{Corollary}
\newcommand{\Fix}{\operatorname{Fix}}

\title{A Coordinate-Sweep Lemma Answers Question 1\\
in ``Fixed Points of Averages of Resolvents''}
\author{Run \texttt{fa\_banach\_001}, lane 8}
\date{July 22, 2026}

\begin{document}
\maketitle

\begin{abstract}
Bauschke, Wang, and Wylie ask whether the fixed-point set of their
Gauss--Seidel-type sweep \(T=J_mR_m\cdots J_1R_1\) equals their
product-space solution set~$S$. We answer the question affirmatively. The
only needed observation is that each map \(J_kR_k\) changes only
coordinate~$k$. More generally, a sweep through maps that update one
distinct coordinate at a time has no fixed points beyond the common fixed
points of its component updates. Combining this elementary lemma with
equation~(64) of the source gives
\[
 S=\Fix T=\Fix(J_mR_m\cdots J_1R_1).
\]
The proof works for every positive integer~$m$ and requires no metric,
linearity, continuity, or monotonicity assumptions on the coordinate
updates.
\end{abstract}

\section{Source question}

The source is H.~H.~Bauschke, X.~Wang, and C.~J.~S.~Wylie,
\emph{Fixed Points of Averages of Resolvents: Geometry and Algorithms},
arXiv:1102.1478v1 \cite{source}. On printed page~19 (PDF page~19), Remark
4.6 asks, for \(m\geq3\):
\begin{quote}
Q1: Concerning (65), is it actually true that
\[
  S=\Fix T=\Fix(J_m\circ R_m\circ\cdots\circ J_1\circ R_1)?
\]
\end{quote}
The source passage is reproduced in Figure~\ref{fig:question}.

\begin{figure}[ht]
\centering
\includegraphics[width=.94\textwidth]{figures/open_problem_crop.png}
\caption{Remark 4.6 and its context in arXiv:1102.1478v1, p.~19.}
\label{fig:question}
\end{figure}

For \(I=\{1,\ldots,m\}\), the source works on the product space
\(\mathbf X=X^m\). Its equations (62)--(63) define maps \(R_k,J_k\) such
that, whenever \(i\neq k\),
\[
  (R_kx)_i=x_i,
  \qquad
  (J_kx)_i=x_i.
\]
Thus \(U_k:=J_k\circ R_k\) changes at most the \(k\)-th coordinate. The
two other source identities needed below are
\begin{equation}\label{eq:sourceS}
 S=\bigcap_{k=1}^m\Fix U_k
 \qquad\text{(source equation (64))}
\end{equation}
and
\begin{equation}\label{eq:sourceT}
 T=U_m\circ\cdots\circ U_1
 \qquad\text{(source equation (66)).}
\end{equation}

\section{Proof intuition}

Run one complete sweep, beginning with coordinate~1. Once coordinate~1
has been updated, no later map can alter it. Therefore, if the end of the
sweep equals the starting vector, the first update must already have left
the vector fixed. After removing that vacuous first update, the same
argument applies to coordinate~2, and successively to all remaining
coordinates. This triangular causality is the entire proof.

\section{The coordinate-sweep theorem}

\begin{theorem}[Fixed points of a one-pass coordinate sweep]
\label{thm:sweep}
Let \(Y=Y_1\times\cdots\times Y_m\) be a Cartesian product of arbitrary
sets. For each \(k\in\{1,\ldots,m\}\), let \(U_k:Y\to Y\) satisfy
\begin{equation}\label{eq:coordinate}
 (U_ky)_i=y_i\qquad(i\neq k).
\end{equation}
Then
\begin{equation}\label{eq:abstractequality}
 \Fix(U_m\circ\cdots\circ U_1)
   =\bigcap_{k=1}^m\Fix U_k.
\end{equation}
\end{theorem}

\begin{proof}
The inclusion from right to left is immediate. For the reverse inclusion,
take \(y\in\Fix(U_m\circ\cdots\circ U_1)\), and define the intermediate
states
\[
 y^{(0)}=y,
 \qquad
 y^{(k)}=U_k y^{(k-1)}\quad(1\leq k\leq m).
\]
Thus \(y^{(m)}=y\). We prove inductively that \(y^{(k)}=y\) for every
\(k\).

The assertion holds for \(k=0\). Suppose \(y^{(k-1)}=y\). By
\eqref{eq:coordinate}, applying \(U_k\) leaves every coordinate other than
\(k\) unchanged, so
\[
 (y^{(k)})_i=y_i\qquad(i\neq k).
\]
Every later map \(U_{k+1},\ldots,U_m\) also leaves coordinate~\(k\)
unchanged. Hence
\[
 (y^{(k)})_k=(y^{(m)})_k=y_k.
\]
All coordinates of \(y^{(k)}\) therefore equal those of~\(y\), proving
\(y^{(k)}=y\). In particular,
\(U_ky=U_ky^{(k-1)}=y^{(k)}=y\). The induction gives
\(y\in\bigcap_k\Fix U_k\), as required.
\end{proof}

\begin{corollary}[Affirmative answer to Question 1]
In the setting of arXiv:1102.1478,
\[
 \boxed{S=\Fix T=\Fix(J_m\circ R_m\circ\cdots\circ J_1\circ R_1).}
\]
In particular, Question 1 in Remark 4.6 has an affirmative answer.
\end{corollary}

\begin{proof}
Set \(U_k=J_k\circ R_k\). Source equations (62)--(63) give
\eqref{eq:coordinate}, so Theorem~\ref{thm:sweep} and
\eqref{eq:sourceT} yield
\[
 \Fix T=\bigcap_{k=1}^m\Fix(J_k\circ R_k).
\]
The right-hand side is \(S\) by \eqref{eq:sourceS}.
\end{proof}

\section{Verification, scope, and novelty}

\paragraph{Definition and order audit.}
The proof uses the source's actual order
\(T=U_m\circ\cdots\circ U_1\): the first applied map is~\(U_1\), and
every later map leaves coordinate~1 unchanged. Equations (62) and (63)
separately leave all non-\(k\) coordinates unchanged, so their composition
does as well. Equation (64) is exactly the common-fixed-point identity
needed in the last line. The argument also remains valid when these sets
are empty.

\paragraph{Strength and limitation.}
The abstract theorem applies to every \(m\geq1\), to any order that visits
each coordinate exactly once, and to maps between products of arbitrary
sets. These strengthenings are incidental. The result answers only Q1 of
Remark 4.6. It supplies no convergence theorem for iterates of~$T$ and
does not answer Q2 or Q3.

\paragraph{Novelty check.}
Cheap run and corpus indexes, followed by bounded searches using the arXiv
identifier, paper title, displayed equality, and core fixed-point and
Gauss--Seidel terminology, found no later general answer. A later paper,
\emph{On moving averages} \cite{moving}, rigorously analyzes the algorithm
in a special case; its abstract describes the result as a special-case
convergence proof. This search is evidence, not an exhaustive
bibliographic guarantee. The result should therefore remain a candidate
full solution pending independent mathematical and literature review.

\begin{thebibliography}{9}
\bibitem{source}
H.~H. Bauschke, X.~Wang, and C.~J.~S. Wylie,
\emph{Fixed Points of Averages of Resolvents: Geometry and Algorithms},
arXiv:1102.1478v1 (2011),
\url{https://arxiv.org/abs/1102.1478}.

\bibitem{moving}
H.~H. Bauschke, J.~Sarada, and X.~Wang,
\emph{On moving averages}, arXiv:1206.3610 (2012),
\url{https://arxiv.org/abs/1206.3610}.
\end{thebibliography}

\end{document}
